\paragraph{倒数矩通道的诱导表示刚性、分拆闭式与逐项卷积公式}
沿用结论链 \ref{subsubsec:xi-exceptional-spectrum-necklace-chain} 的记号。固定 \(q\ge 2\)，记
\[
C^{(q)}_{i,j}:=\binom{q-i}{j}\qquad(0\le i,j\le q),\qquad
A^{(q)}=\frac12\Bigl(\mathbf 1(b^{(q)})^{\mathsf T}+C^{(q)}\Bigr),
\]
并定义
\[
U^{(q)}:=2\bigl(C^{(q)}\bigr)^{-1},\qquad
S_{-m}(q):=\Tr\!\bigl((A^{(q)})^{-m}\bigr)\quad(m\ge1).
\]
采用延拓 Fibonacci 约定
\[
F_{-1}:=1,\qquad F_0:=0,\qquad F_1:=1.
\]
此外定义倒数 \(\Omega\)-通道
\[
\Omega_m^{(-)}(q):=\Tr\!\bigl((C^{(q)})^{-m}\bigr),
\]
则由 \(\mathrm{spec}(C^{(q)})=\{(-1)^j\varphi^{q-2j}\}_{j=0}^{q}\) 可得
\[
\Omega_m^{(-)}(q)=(-1)^{mq}\Omega_m(q),
\]
其中 \(\Omega_m(q)\) 见结论 \ref{con:xi-terminal-replica-softcore-lucas-pell-omega-unified}。

\begin{theorem}[诱导表示角元刚性与 Fibonacci 指数律]\label{thm:xi-replica-softcore-reciprocal-corner-fibonacci-power-law}
对任意整数 \(m\ge0\)，成立严格恒等式
\[
\boxed{\;
\bigl(U^{(q)}\bigr)^m_{0,0}
=2^m(-1)^{mq}F_{m-1}^{\,q}\; }.
\]
\end{theorem}
\leanverified{paper\_xi\_replica\_softcore\_reciprocal\_corner\_fibonacci\_power\_law}
\begin{proof}
令
\[
Q:=\begin{pmatrix}1&1\\[2pt]1&0\end{pmatrix}.
\]
在齐次多项式空间 \(\Sym^q(\CC^2)\) 的基
\(
\{x^{q-j}y^j\}_{j=0}^{q}
\)
下，二项展开给出
\[
\Sym^q(Q)=\bigl(C^{(q)}\bigr)^{\mathsf T},
\]
从而
\[
\bigl(C^{(q)}\bigr)^{-1}
=\bigl(\Sym^q(Q^{-1})\bigr)^{\mathsf T}.
\]
于是
\[
\bigl(U^{(q)}\bigr)^m
=2^m\Bigl(\Sym^q(Q^{-1})^m\Bigr)^{\mathsf T}
=2^m\Bigl(\Sym^q(Q^{-m})\Bigr)^{\mathsf T}.
\]
另一方面，Fibonacci 矩阵恒等式给出
\[
Q^m=
\begin{pmatrix}
F_{m+1}&F_m\\
F_m&F_{m-1}
\end{pmatrix},
\qquad
Q^{-m}=(-1)^m
\begin{pmatrix}
F_{m-1}&-F_m\\
-F_m&F_{m+1}
\end{pmatrix}.
\]
设 \(Q^{-m}=\begin{psmallmatrix}\alpha&\beta\\ \gamma&\delta\end{psmallmatrix}\)，则
\(
\bigl(\Sym^q(Q^{-m})\bigr)_{0,0}
\)
等于 \((\alpha x+\beta y)^q\) 的 \(x^q\) 系数，即 \(\alpha^q\)。
由 \(\alpha=(-1)^mF_{m-1}\)，得到
\[
\bigl(U^{(q)}\bigr)^m_{0,0}
=2^m\alpha^q
=2^m(-1)^{mq}F_{m-1}^{\,q}.
\]
证毕。
\end{proof}

\begin{theorem}[例外谱倒数矩的 Fibonacci--partition 分拆闭式]\label{thm:xi-replica-softcore-reciprocal-moment-partition-closed-form}
\leanverified{paper\_xi\_replica\_softcore\_reciprocal\_moment\_partition\_closed\_form}
固定 \(m\ge1\)。对分拆
\[
\pi=1^{m_1}2^{m_2}\cdots m^{m_m}\vdash m,
\]
记
\[
\ell(\pi):=\sum_{s=1}^{m}m_s,\qquad
c(\pi):=\frac{m}{\ell(\pi)}\cdot\frac{\ell(\pi)!}{\prod_{s=1}^{m}m_s!},
\]
\[
\Psi(\pi):=\prod_{s=1}^{m}F_{s-2}^{\,m_s}.
\]
则对任意 \(q\ge2\) 成立
\[
\boxed{\;
S_{-m}(q)
=2^m\Omega_m^{(-)}(q)
+\sum_{\substack{\pi\vdash m\\ \ell(\pi)\ge1}}
(-1)^{\ell(\pi)+q(m-\ell(\pi))}
c(\pi)\,2^{m-\ell(\pi)}\,\Psi(\pi)^q\; }.
\]
\end{theorem}
\begin{proof}
由定理 \ref{thm:xi-replica-softcore-exceptional-block-inverse-integral-closed}，
\[
\bigl(A^{(q)}\bigr)^{-1}=U^{(q)}-E,\qquad E:=E_{00}=e_0e_0^{\mathsf T},
\]
故
\[
\bigl(A^{(q)}\bigr)^{-m}=(U^{(q)}-E)^m.
\]
对该式作非交换词展开并取迹。

不含 \(E\) 的唯一词为 \(U^{(q)m}\)，贡献
\[
\Tr\!\bigl(U^{(q)m}\bigr)
=2^m\Tr\!\bigl((C^{(q)})^{-m}\bigr)
=2^m\Omega_m^{(-)}(q).
\]

设词中 \(E\) 出现 \(\ell\ge1\) 次。利用迹的循环不变性，可把每个循环类唯一写成块长
\(
(s_1,\dots,s_\ell)
\)
（\(s_r\ge1\)、\(\sum_r s_r=m\)），其意义是每块为 \(E(U^{(q)})^{s_r-1}\)。
记对应分拆为 \(\pi\vdash m\)。
由
\[
E(U^{(q)})^aE=\bigl((U^{(q)})^a\bigr)_{0,0}E
\]
和定理 \ref{thm:xi-replica-softcore-reciprocal-corner-fibonacci-power-law}，
\[
\Tr\!\bigl(E(U^{(q)})^{s_1-1}\cdots E(U^{(q)})^{s_\ell-1}\bigr)
=\prod_{r=1}^{\ell}2^{s_r-1}(-1)^{q(s_r-1)}F_{s_r-2}^{\,q}
=2^{m-\ell}(-1)^{q(m-\ell)}\Psi(\pi)^q.
\]
再乘上 \((U^{(q)}-E)^m\) 中 \(\ell\) 次选取 \((-E)\) 的符号 \((-1)^\ell\)，得到每个循环类贡献
\[
(-1)^\ell\,2^{m-\ell}(-1)^{q(m-\ell)}\Psi(\pi)^q.
\]
固定分拆 \(\pi\) 时，其循环类计数为
\[
c(\pi)=\frac{m}{\ell(\pi)}\cdot\frac{\ell(\pi)!}{\prod_s m_s!}.
\]
对全部分拆求和即得结论。
\end{proof}

\begin{theorem}[例外谱倒数幂和的环分解公式]\label{thm:xi-replica-softcore-reciprocal-moment-ring-decomposition}
\leanverified{paper\_xi\_replica\_softcore\_reciprocal\_moment\_ring\_decomposition}
固定 \(m\ge 1\)。令
\[
s_n:=(-1)^{nq}F_{n-1}^{\,q}\qquad(n\ge 0),
\]
并约定 \(F_{-1}=1\)。
则对任意 \(q\ge 2\) 成立等价闭式
\[
\boxed{\;
S_{-m}(q)
=2^m\Omega_m^{(-)}(q)
\sum_{k=1}^{m}(-1)^k\,2^{m-k}\,\frac{m}{k}
\sum_{\substack{n_1+\cdots+n_k=m-k\\ n_j\ge 0}}
\ \prod_{j=1}^{k}s_{n_j}\; }.
\]
特别地，\(S_{-m}(q)\in\ZZ\) 对一切 \(m\ge 1\)、\(q\ge 2\) 成立。
\end{theorem}
\begin{proof}
沿用定理 \ref{thm:xi-replica-softcore-reciprocal-moment-partition-closed-form} 证明中的展开
\(
(A^{(q)})^{-m}=(U^{(q)}-E)^m
\)
与循环类归并，但不再按分拆 \(\pi\) 进一步合并。

\par
设在某一项中 \(E\) 出现 \(k\ge 1\) 次。
由迹的循环不变性，可把该循环类唯一编码为一个 \(k\)-元非负整数组合 \((n_1,\dots,n_k)\) 满足 \(\sum_{j=1}^{k}n_j=m-k\)，其意义为相邻两次 \(E\) 之间出现 \(U^{(q)}\) 的次数。
每个循环类在长度为 \(m\) 的线性展开中出现次数为 \(m/k\)。
又由
\[
E\bigl(U^{(q)}\bigr)^{n}E=\bigl(U^{(q)}\bigr)^{n}_{0,0}\,E
=2^n(-1)^{nq}F_{n-1}^{\,q}\,E
\]
（定理 \ref{thm:xi-replica-softcore-reciprocal-corner-fibonacci-power-law}，并用延拓约定使 \(n=0\) 时一致），得到该循环类对迹的贡献为
\[
(-1)^k\,2^{m-k}\prod_{j=1}^{k}\bigl((-1)^{n_jq}F_{n_j-1}^{\,q}\bigr),
\]
其中 \((-1)^k\) 来自选取 \((-E)\) 的符号。
对所有 \(k\) 与 \((n_1,\dots,n_k)\) 求和并加上无 \(E\) 的项 \(2^m\Omega_m^{(-)}(q)\)，即得所述公式。
整数性由 \( (A^{(q)})^{-1}\in\ZZ^{(q+1)\times(q+1)} \) 推出：
\(
S_{-m}(q)=\Tr(((A^{(q)})^{-1})^m)\in\ZZ
\)。
\end{proof}

\begin{corollary}[倒数三至六次矩的显式闭式]\label{cor:xi-replica-softcore-reciprocal-moment-m3-m6-closed}
\leanverified{paper\_xi\_replica\_softcore\_reciprocal\_moment\_m3\_m6\_closed}
对任意 \(q\ge2\)，成立
\[
\boxed{\;
S_{-3}(q)=8\,\Omega_3^{(-)}(q)-13\;},
\]
\[
\boxed{\;
S_{-4}(q)=16\,\Omega_4^{(-)}(q)-32(-1)^q+17\;},
\]
\[
\boxed{\;
S_{-5}(q)=32\,\Omega_5^{(-)}(q)-80\cdot 2^q+40(-1)^q-21\;},
\]
\[
\boxed{\;
S_{-6}(q)=64\,\Omega_6^{(-)}(q)+96\cdot 2^q-192(-1)^q3^q-48(-1)^q+73\; }.
\]
\end{corollary}
\begin{proof}
在定理 \ref{thm:xi-replica-softcore-reciprocal-moment-partition-closed-form} 中分别取 \(m=3,4,5,6\)。
含部件 \(2\) 的分拆项因 \(F_0=0\) 全部消失；其余有限项直接化简即得。
\end{proof}

\begin{corollary}[倒数一次至八次矩的完全闭式]\label{cor:xi-replica-softcore-reciprocal-moment-m1-m8-closed}
对任意 \(q\ge 2\)，成立
\[
\boxed{\;
S_{-1}(q)=2\,\Omega_1^{(-)}(q)-1=2(-1)^qF_{q+1}-1\;},
\]
\[
\boxed{\;
S_{-2}(q)=4\,\Omega_2^{(-)}(q)+1=4F_{2q+2}+1\;},
\]
以及推论 \ref{cor:xi-replica-softcore-reciprocal-moment-m3-m6-closed} 的 \(m=3,4,5,6\) 公式；并且
\[
\boxed{\;
S_{-7}(q)=128\,\Omega_7^{(-)}(q)-448\cdot 5^{q}+224(-1)^q 3^{q}-112\cdot 2^{q}+280(-1)^q-141\;},
\]
\[
\boxed{\;
S_{-8}(q)=256\,\Omega_8^{(-)}(q)-1024(-1)^q 8^{q}+512\cdot 5^{q}-256(-1)^q 3^{q}+640\cdot 2^{q}-576(-1)^q+481\; }.
\]
\end{corollary}
\begin{proof}
在定理 \ref{thm:xi-replica-softcore-reciprocal-moment-partition-closed-form} 中分别取 \(m=1,2,7,8\)。
当 \(m=1\) 时仅保留 \(\ell(\pi)=1\) 的常数项，得到 \(S_{-1}(q)=2\Omega_1^{(-)}(q)-1\)；
当 \(m=2\) 时仅出现 \(\Psi(\pi)\in\{F_{-1},F_0\}\) 且 \(F_0=0\) 使非平凡分拆项消失，得到 \(S_{-2}(q)=4\Omega_2^{(-)}(q)+1\)。
当 \(m=7,8\) 时同理由有限分拆项直接化简得到所示闭式。
\end{proof}

\begin{theorem}[全矩阵逆幂迹的 Lucas 主项与有限修正]\label{thm:xi-replica-softcore-full-trace-negative-power-lucas-finite-correction}
令 \(q\ge 2\)，并记 Lucas 数列 \(L_0=2\)、\(L_1=1\)、\(L_{m+1}=L_m+L_{m-1}\)。
则对任意整数 \(m\ge 1\) 成立
\[
\boxed{\;
\Tr\!\bigl((T^{(q)})^{-m}\bigr)
=2^m(-1)^{mq}L_m^{\,q}
\sum_{k=1}^{m}(-1)^k\,2^{m-k}\,\frac{m}{k}
\sum_{\substack{n_1+\cdots+n_k=m-k\\ n_j\ge 0}}
\ \prod_{j=1}^{k}\Bigl((-1)^{n_jq}F_{n_j-1}^{\,q}\Bigr)\; }.
\]
从而对 \(m=1,\dots,8\)（其中 \(L_1=1,L_2=3,L_3=4,L_4=7,L_5=11,L_6=18,L_7=29,L_8=47\)）有完全闭式
\[
\begin{aligned}
\Tr\!\bigl((T^{(q)})^{-1}\bigr)&=2(-1)^q-1,\\
\Tr\!\bigl((T^{(q)})^{-2}\bigr)&=4\cdot 3^q+1,\\
\Tr\!\bigl((T^{(q)})^{-3}\bigr)&=8(-1)^q\cdot 4^q-13,\\
\Tr\!\bigl((T^{(q)})^{-4}\bigr)&=16\cdot 7^q+17-32(-1)^q,\\
\Tr\!\bigl((T^{(q)})^{-5}\bigr)&=32(-1)^q\cdot 11^q-80\cdot 2^q-21+40(-1)^q,\\
\Tr\!\bigl((T^{(q)})^{-6}\bigr)&=64\cdot 18^q+96\cdot 2^q-192(-1)^q 3^q+73-48(-1)^q,\\
\Tr\!\bigl((T^{(q)})^{-7}\bigr)&=128(-1)^q\cdot 29^q-448\cdot 5^q+224(-1)^q 3^q-112\cdot 2^q+280(-1)^q-141,\\
\Tr\!\bigl((T^{(q)})^{-8}\bigr)&=256\cdot 47^q-1024(-1)^q 8^q+512\cdot 5^q-256(-1)^q 3^q+640\cdot 2^q-576(-1)^q+481.
\end{aligned}
\]
\end{theorem}
\begin{proof}
由定理 \ref{thm:xi-replica-softcore-full-spectrum-direct-sum} 的例外--体谱直和分解，
\[
\Tr\!\bigl((T^{(q)})^{-m}\bigr)
=S_{-m}(q)+\sum_{i=0}^{q}\Bigl(\binom{q}{i}-1\Bigr)\Bigl(\frac{2}{\mu_i^{(q)}}\Bigr)^{m},
\qquad
\mu_i^{(q)}=\varphi^{\,q-i}\psi^{\,i}.
\]
体谱部分满足
\[
\sum_{i=0}^{q}\binom{q}{i}\bigl(\mu_i^{(q)}\bigr)^{-m}
=(\varphi^{-m}+\psi^{-m})^q
=\bigl((-1)^mL_m\bigr)^q
=(-1)^{mq}L_m^{\,q},
\]
而
\(
\sum_{i=0}^{q}\bigl(\mu_i^{(q)}\bigr)^{-m}
=\Tr\!\bigl((C^{(q)})^{-m}\bigr)=\Omega_m^{(-)}(q)
\)
（见本段定义）。
因此体谱贡献为
\(
2^m\bigl((-1)^{mq}L_m^{\,q}-\Omega_m^{(-)}(q)\bigr)
\)。
将其与定理 \ref{thm:xi-replica-softcore-reciprocal-moment-ring-decomposition} 中的 \(S_{-m}(q)\) 相加，\(\Omega_m^{(-)}(q)\) 项严格相消，得到一般公式；将 \(m=1,\dots,8\) 的有限和逐一化简即可得到列出的闭式。证毕。
\end{proof}

\begin{theorem}[固定阶倒数矩的 \(q\)-方向递推因子]\label{thm:xi-replica-softcore-reciprocal-moment-q-recurrence}
\leanverified{paper\_xi\_replica\_softcore\_reciprocal\_moment\_q\_recurrence}
固定 \(m\ge1\)。定义
\[
\mathcal A_m^{-}:=
\Bigl\{
(-1)^{m-\ell(\pi)}\Psi(\pi)\,:\,
\pi\vdash m,\ \ell(\pi)\ge1,\ \Psi(\pi)\neq0
\Bigr\}\subset\ZZ.
\]
记 Lucas 数 \(L_m=\varphi^m+(-\varphi^{-1})^m\)，并置
\[
\widetilde L_m:=(-1)^mL_m.
\]
则序列 \(q\mapsto S_{-m}(q)\) 满足常系数线性递推，其特征多项式可取为
\[
\boxed{\;
P_m^{-}(X)=\bigl(X^2-\widetilde L_m X+(-1)^m\bigr)
\prod_{\alpha\in\mathcal A_m^{-}}(X-\alpha)\; }.
\]
将重复根合并并去除零系数通道后，得到该序列的最小特征多项式。
\end{theorem}
\begin{proof}
由定理 \ref{thm:xi-replica-softcore-reciprocal-moment-partition-closed-form}，
\[
S_{-m}(q)=2^m\Omega_m^{(-)}(q)+\sum_{\alpha\in\mathcal A_m^{-}}B_\alpha\,\alpha^q,
\]
其中 \(B_\alpha\in\QQ\) 与 \(q\) 无关。
又
\[
\Omega_m^{(-)}(q)
=u_m(\varphi^{-m})^q+v_m\bigl(((-1)^m\varphi^{m})\bigr)^q
\]
（\(u_m,v_m\neq0\)），故其特征二次因子为
\[
(X-\varphi^{-m})(X-(-1)^m\varphi^m)
=X^2-\widetilde L_m X+(-1)^m.
\]
每个指数项 \(\alpha^q\) 被一次因子 \(X-\alpha\) 湮灭，
因此 \(P_m^{-}\) 湮灭整个序列。合并重复根并删除零系数通道后即得最小特征多项式。
\end{proof}

\begin{theorem}[逆二项矩阵负幂的逐项卷积闭式]\label{thm:xi-replica-softcore-binomial-inverse-negative-power-entrywise}
对任意 \(q\ge2\)、\(m\ge0\) 与 \(0\le i,j\le q\)，成立
\[
\boxed{
\begin{aligned}
\bigl(C^{(q)}\bigr)^{-m}_{i,j}
=(-1)^{mq+i+j}
\sum_{t=\max(0,i+j-q)}^{\min(i,j)}
&\binom{i}{t}\binom{q-i}{j-t}\\
&\cdot F_{m-1}^{\,q-i-j+t}
F_m^{\,i+j-2t}
F_{m+1}^{\,t}.
\end{aligned}
}
\]
等价地，
\[
\bigl(U^{(q)}\bigr)^m_{i,j}=2^m\bigl(C^{(q)}\bigr)^{-m}_{i,j}.
\]
\end{theorem}
\leanverified{paper\_xi\_replica\_softcore\_binomial\_inverse\_negative\_power\_entrywise}
\begin{proof}
仍记
\(
Q=\begin{psmallmatrix}1&1\\1&0\end{psmallmatrix}
\)。
由 \(\Sym^q(Q)=\bigl(C^{(q)}\bigr)^{\mathsf T}\)，得
\[
\bigl(C^{(q)}\bigr)^{-m}
=\Bigl(\Sym^q(Q^{-m})\Bigr)^{\mathsf T}.
\]
设
\[
Q^{-m}=
\begin{pmatrix}
a&b\\ c&d
\end{pmatrix},
\]
则 \(\bigl(\Sym^q(Q^{-m})\bigr)_{j,i}\) 等于
\(
(ax+by)^{q-i}(cx+dy)^i
\)
中 \(x^{q-j}y^j\) 的系数，故
\[
\bigl(C^{(q)}\bigr)^{-m}_{i,j}
=\sum_{t}
\binom{i}{t}\binom{q-i}{j-t}
a^{\,q-i-(j-t)}b^{\,j-t}c^{\,i-t}d^{\,t},
\]
其中求和区间由指数非负性给出
\(
t\in[\max(0,i+j-q),\min(i,j)].
\)
代入
\[
Q^{-m}=(-1)^m
\begin{pmatrix}
F_{m-1}&-F_m\\ -F_m&F_{m+1}
\end{pmatrix},
\]
并整理符号指数，即得
\(
(-1)^{mq+i+j}
\)
与所示三参数卷积因子。证毕。
\end{proof}

\endinput
